\section{Introduction}
\label{sec:intro}

The modern paradigm of large-scale deep learning serving is increasingly shifting toward modular, multi-expert architectures. By decomposing a monolithic neural network into specialized parameter-efficient experts (e.g., via Low-Rank Adaptation (LoRA) \cite{hu2021lora} or Mixture-of-Experts (MoE) \cite{shazeer2017outrageously, fedus2022switch}), machine learning systems can serve highly diverse and non-stationary query streams without the catastrophic interference or massive memory overhead associated with running multiple independent models. In these architectures, a dynamic routing layer plays a central, orchestrating role: it extracts intermediate activations from a shared backbone, projects them onto task-representative coordinate spaces, and dynamically computes a set of ensembling weights to blend the expert activations at runtime \cite{wortsman2022model, ilharco2022editing}.

However, operating such adaptive serving layers in real-world sequential streaming environments introduces an acute, fundamental challenge known as the \emph{Jitter-Lag Trade-Off}. Standard stateless routers, such as SABLE \cite{sable2024}, evaluate each incoming query in absolute isolation. By mapping instantaneous sensory projections to gating coefficients via direct nearest-centroid or softmax similarity metrics, stateless designs react immediately to task boundaries. Yet, this high sensitivity is a double-edged sword: minor representation perturbations, out-of-distribution observation noise, and high-frequency activation fluctuations propagate directly to the gating weights. This results in severe \emph{routing jitter}---rapid, high-frequency oscillations of the ensembling coefficients across adjacent steps in a sequence. Beyond its impact on isolated classification, this chaotic switching is hypothesized to introduce severe systems-level and representational bottlenecks in physical settings: (1) \emph{Potential Hardware Cache Thrashing:} Under dynamic Mixture-of-Experts (MoE) or parameter-efficient serving (e.g., loading LoRA weights dynamically), high-frequency weight oscillations are expected to trigger continuous reloading, scaling, or swapping of adapter parameters in high-speed GPU SRAM or L1 cache, potentially degrading memory coalescing and parallel serving efficiency. (2) \emph{Hypothesized Representational Instability:} In deep, multi-layer architectures, rapid layer-by-layer ensembling weight fluctuations may disrupt the distribution of intermediate activations, potentially breaking the semantic and functional coherence of representations across network depth---an effect that static, single-step evaluations do not capture but represents a critical risk for deep network stability \cite{packinetics2025}.

To mitigate routing noise, recent research has turned to stateful gating methods. Inspired by continuous-time physical systems, frameworks like ChemMerge \cite{chemmerge2025} model expert concentrations as chemical species governed by non-equilibrium reaction kinetics and ordinary differential equations (ODEs). Under mathematical simplification, such complex ODEs can be deconstructed to a constant exponential moving average (EMA) of gating weights across sequential depth, a paradigm formalized as Momentum-Merge \cite{momentummerge2025}. While these temporal filters successfully act as low-pass filters to smooth ensembling trajectories, they do so at a devastating cost: \emph{representational lag (inertial drag)}. Because stateful temporal filters accumulate history rigidly, they are incapable of distinguishing high-frequency noise from a genuine, abrupt task transition. Consequently, when a serving stream undergoes a sharp context switch, stateful routers exhibit a massive delay in updating their active mixture coefficients, leading to a catastrophic collapse in accuracy as the system continues to process queries under the stale representation of the previous task.

In this work, we argue that this trade-off is not an insurmountable physical law, but rather an artifact of a fundamentally passive routing paradigm. Existing routers are designed either as feedforward mathematical mappings or as rigid, unyielding temporal smoothers. In strict alignment with a visionary philosophy that seeks to rethink core assumptions, we propose a radical departure from these passive heuristics. We argue that the multi-expert routing layer should instead be modeled as an \emph{active, self-organizing cognitive agent} performing test-time perception and action in a dynamic, non-i.i.d. environment.

To realize this vision, we introduce \textbf{Active Inference Routing (AIR)}, a novel brain-inspired framework grounded in the Free Energy Principle from theoretical neuroscience and cognitive science \cite{friston2010free, friston2006free}. Rather than passing activations through isolated feedforward layers, AIR formalizes the serving stream as a state-space dynamical system where task contexts are latent hidden states to be perceived, and expert gating configurations represent the actions taken by the routing agent. 
The agent maintains a stateful, Gaussian variational belief state tracking the latent log-probabilities of experts. Upon receiving a noisy sensory activation, the agent computes the \textbf{Variational Free Energy}---a principled informational metric that mathematically balances top-down prior transition expectations with bottom-up sensory coordinate projections. By solving this convex free energy objective exactly in a single closed-form step, the routing agent performs active perception, updating its belief state in a highly adaptive, context-dependent manner.

This active minimization elegantly resolves the jitter-lag dilemma. During stationary periods, when prediction errors are low, top-down prior expectations dominate, filtering out high-frequency sensory noise and guaranteeing smooth, stable ensembling weights. However, upon encountering a genuine task boundary, the bottom-up sensory coordinate predictions diverge drastically from prior expectations, causing a massive, instantaneous spike in the prediction error. This prediction error acts as an informational force that immediately overcomes the temporal prior, resetting the belief state and driving near-instantaneous adaptation (within 1--2 steps) without any representational lag.

\textbf{Summary of Contributions:}
\begin{itemize}
    \item \textbf{Brain-Inspired Gating Framework:} We formulate the first multi-expert serving routing layer as an active-inference cognitive agent, redefining model merging as active perception.
    \item \textbf{Variational Free Energy Formulation:} We analytically derive the Variational Free Energy for serving streams, showing how it simplifies to precision-weighted squared prediction errors that balance stability and adaptation.
    \item \textbf{Resolution of the Jitter-Lag Trade-Off:} On the Analytical Coordinate Sandbox (ACS), we show that AIR matches oracle representation alignment under rapid switches, while slashing SABLE's routing jitter by up to 2.49$\times$ under stable, noisy streams.
    \item \textbf{Verification of Inhibitory Pathways:} We validate the mechanistic necessity of active inhibition in the generative mapping, proving that restricting the mapping to non-negative configurations introduces localized transient lag at switch boundaries.
\end{itemize}
